\chapter{Conclusion}
\label{ch:conclusion}

\section{Conclusions and contributions}

Method selection should precede serialization optimization in LLM-based information extraction from
structured data. A compact representation can make a direct input more practical, but the application's
tasks and constraints determine whether the extraction method is usable.

We contributed a controlled comparison of Direct Context Inference and Text-to-SQL and developed the
evidence-informed decision framework in Section~\ref{sec:decision-framework}. Within the scope discussed
in Chapter~\ref{ch:discussion}, Text-to-SQL is the preferred starting point for relational
workloads where the tasks and infrastructure support SQL. For feasible Direct Context Inference over flat
data, the findings favour compact representations, with CSV as a starting choice and TOON as a competitive
alternative. This is conditional guidance, not a universal ordering of methods or formats.

The framework makes the assumptions behind a starting choice explicit. It identifies candidates for
evaluation, not a configuration that satisfies every application. Validation on representative questions
is still required before adoption.

\section{Future work}

The issues discussed in Section~\ref{sec:validity-limitations} motivate three directions for further work.

First, future experiments should use relationship-preserving sampling or full-scale databases to examine
whether the conclusions hold when joins and aggregates operate on intact relationships. Comparison with
independently sampled tables would help isolate the effects of sampling. Answer evaluation should also
compare structured results directly, preserving
numeric values and handling row order and duplicates according to the question's semantics. Applying this
comparison to both pipelines would clarify how much of the observed agreement difference reflects
correct retrieval rather than differences in answer representation.

Second, the pipelines should be extended with selective retrieval, targeted metadata, and semantic
verification. Controlled additions would show whether supplying relevant records or field descriptions,
or checking the meaning of a generated query, addresses errors left unresolved by the baseline.
Evaluation should account for the additional preprocessing and inference work, including how preprocessing
is shared across repeated requests. Nested-data retrieval should also be tested with structure-preserving
formats, including TOON, to establish whether the framework's conditional format branch provides useful
guidance beyond flat relational inputs. These experiments would clarify when added complexity is justified.

Third, the framework should be evaluated on independent application workloads by testing whether its
suggested candidates meet the stated requirements. Accuracy and resource use should be measured on the
same question set so that application trade-offs can be assessed directly. Realistic serving conditions,
including concurrent requests and repeated access to data, would test whether the guidance transfers
to deployment and identify selection rules that need revision.

The thesis provides an evidence-informed starting point for choosing how LLMs extract information from
structured data.
